\paragraph{The rank removed by the arithmetic projector.}

A separating section identifies the levels. The projector can still remove part of their rank, which is counted in the finite representation at the two level primes.

Use the finite group of the two level primes in the same model.
For this calculation twist its action by the inverse of the
central form character $\chi$, so that the chosen character
becomes the product Whittaker character $\lambda$ of $U$. At
each residue factor its value is a nontrivial additive character
of the sum on the first superdiagonal. This twist only multiplies
the scalar in the pullback action. Write $V$ for the cusp space
of weight $\kappa$ with that action and $V_\lambda$ for its
character space. Its dimension is $d_0$.

\begin{lemma}[Trace of the discarded central idempotent]
\label{lem:paired-projector-rank-loss}
For the character space of dimension \(d_0\), write \(r_{\rm bad}\) for the rank discarded by the bad-type projector. The ratios \(\beta_k\) from \eqref{eq:row-bad-share} give
\[
 r_{\rm bad}\le2d_0(\beta_0+\beta_1).
\]
 This comparison uses the integral
arithmetic projector through its scalar extension.
\end{lemma}

\begin{proof}
\proofstep{The finite type decomposition at the two level primes.}
Let $H_k$ be the whole finite linear group at the $k$th level
prime and the upper unitriangulars at the other. There is an
$H_k$-equivariant embedding of $V$ into $a_k$ copies of the
regular representation, with $a_k |H_k|\le 2|U|d_0$.
For a definite form space use the free action on components,
which gives equality with $|U|d_0$ in this bound. In the
indefinite case, let $w$ be a separating section for $H=H_k$.
Write $w_h$ for its ordinary pullbacks, indexed so that the
ring action of $g$ sends $w_h$ to $w_{gh}$. The scalar twist belongs to the action on the form \(v\), while pullback permutes the separating sections. Consequently
\[
  g(w_h v)=w_{gh}\,g(v).
\]

Let $W$ be the invariant subspace of the span of all $w_h v$
for $v\in V$. Consider the map
\[
 V\longrightarrow \operatorname{Map}(H,W),\qquad
 v\longmapsto
    \left(\sum_{g\in H}g(w_h v)\right)_{h\in H}.
\]
For $kv$ the tuple at $h$ becomes the tuple for $v$ at
$k^{-1}h$, so the map intertwines the left regular action.
If the tuple vanishes, the group determinant of the $w_h$
gives $g(v)=0$ for each $g$ on a dense open set of every
component. This gives injectivity. Conversely the map
$(a_h)\mapsto\sum_h w_h a_h$ from $W^{\oplus H}$ has target weight $\kappa+2d$.

If its value vanishes, applying each $g\in H$ gives the same
invertible matrix of translates on the $a_h$ because $W$
consists of invariants. Writing \(a_k=\dim W\), apply first the injection just constructed,
then the three-weight comparison and finally the trace injection for \(U\):
\[
 a_k|H_k|\le\dim V_{\kappa+2d}
 \le2\dim V_{\kappa-d}\le2|U|d_0 .
\]
The bounds on the irreducible multiplicities for $H_k$ are
unchanged by extending the algebraic coefficient field through
a different complex embedding. Indeed, over an algebraic
closure they are the ranks of the central primitive idempotents
divided by the corresponding irreducible degrees. A change of
embedding permutes these idempotents and preserves their degrees.
The regular representation contains each irreducible with
multiplicity its degree. Thus the embedding bound also holds
at the embedding used for the finite type projectors.
\label{cor:paired-good-share-1}

\proofstep{Tracing the central bad projector.}
In the group algebra let $b_k$ be the central idempotent for
bad tuples of semisimple labels at the $k$th prime, extended
by the identity over the other prime. The two commuting central idempotents give the retained projector
\[
  e_{\rm good}=(1-b_0)(1-b_1)=1-b_0-b_1+b_0b_1.
\]
 Let $e_\lambda$
be the normalized character idempotent of $U$.
The operator $b_ke_\lambda$ is idempotent and belongs to
$\mathbb C[H_k]$. The dimension discarded inside $V_\lambda$
is at most the sum of the two ranks of these idempotents.

For the regular representation of $H_k$ the trace of group
algebra multiplication by $B$ is $|H_k|$ times the coefficient
of $1$ in $B$. Indeed each diagonal entry on the group basis
is that coefficient. Consequently, writing $\beta_k$ for that
coefficient in $b_ke_\lambda$ divided by the coefficient in
$e_\lambda$, the rank discarded is at most
\begin{equation}\label{eq:row-bad-share}
 \sum_k a_k\frac{|H_k|}{|U|}\,\beta_k
 \ \le\ 2d_0(\beta_0+\beta_1).
\end{equation}
Here the projector image has its arithmetic meaning even if the
cusp reading uses a different embedding from the type weights.
Both the character kernel and its intersection with the kernel of
the good idempotent are ranks of stacked algebraic linear maps.
This is the use of the two embeddings in the selected comparison.
\label{cor:paired-good-share-2}

\end{proof}

Let \(G=\operatorname{GL}_r(\mathbb F_g)\), \(g=g_k\), with upper
unitriangular subgroup \(U\). For a monic degree-\(r\) polynomial
\(P\) with unit constant coefficient let \(e_P\) be its central
series projector. Write \(e_\psi\) for the idempotent of a standard
nondegenerate character of \(U\) and \([1]\) for the coefficient
of the identity in a group-algebra element.

\begin{lemma}[Local weights of squarefree type labels]
The numbers \(\omega_g(P)=[1](e_Pe_\psi)/[1]e_\psi\) are
nonnegative, sum to one and are independent of the standard
nontrivial additive character used to define \(\psi\). If \(P\)
is squarefree, then
\[
        \omega_g(P)=\prod_{Q\mid P}(g^{\deg Q}-1)^{-1},
\]
where \(Q\) ranges over monic irreducible factors.
\end{lemma}
\begin{proof}
In the regular representation \(\operatorname{tr}_{\rm reg}(a)
=|G|[1]a\); the projectors \(e_Pe_\psi\) split the range of
\(e_\psi\). Conjugation by a diagonal matrix with the prescribed
successive ratios carries one standard character to the other,
fixing each \(e_P\). This proves the first assertions.

For squarefree \(P\) the dual element has distinct eigenvalues;
its centralizer is the torus \(T=\prod_{Q\mid P}\mathbb F_{g^{\deg Q}}^\times\).
The torus character is in general position. The Deligne--Lusztig
degree and Gelfand--Graev multiplicity formulas for \(G\), whose
center is connected, give one irreducible in this series of degree
\(|G|_{g'}/|T|\), with Whittaker multiplicity one. Thus
\[
 \omega_g(P)=\frac{(|G|_{g'}/|T|)\cdot1}{|G|/|U|}=|T|^{-1}.
\]
\cite[\S\S5.24, 7.1, 7.4 and 10.7]{deligne-lusztig-finite}
\end{proof}
At the level prime \(g_k\), sum the products of these local
weights over bad tuples, one factor per split place, to obtain
\(\beta_k\). \label{cor:paired-good-share-3}

\paragraph{Bounds for the four finite label exclusions.}
Lemma~\ref{lem:finite-type-event-bounds} gives the four mass
bounds used here.
The finite type bound used on the selected family is
\begin{align}
 \beta_k\le{}&\frac{3f}{g}+\frac{23 f\log g}{\sqrt g}
     +2fm\left(\frac1{4\min(q,d')}+\frac{1+\log r}{g-1}\right)
       \notag\\
     &+\frac r2(7/8)^f+\exp\left(-\frac{f}{5\log r}\right).
 \label{eq:row-bad-types}
\end{align}
In this formula $q,d'$ are the fixed auxiliary primes at the
point, and $d'$ is distinct from the determinant weight $d$.
The hypotheses are $r\ge100$, $g\ge4096$, distinct odd $q,d'$
with $\max(q,d')\ge16$ and $16qd'\mid g-1$, and the same
character labels as in the projector. The squarefree mass is used
as a subprobability. In degree $r\ge2$ there are $g^r-g^{r-1}$
monic squarefree polynomials by the Euler product, and removing
zero constant terms costs at most $g^{r-1}$.

Since the torus
order of each remaining label is at most $g^r$, the missing
weight is at most $2/g$.
The retained mass is consequently bounded below by
\[
  \frac{g^r-g^{r-1}-g^{r-1}}{g^r}=1-\frac2g .
\]
 A union over the $f$ places fits the
first displayed term. The other terms use four events on the
squarefree labels, namely a pair of roots with equal prime-to-$\ell$
parts at a nonexcluded prime $\ell$, a matching part against a
prescribed scalar or its inverse, a common proper divisor degree
at all places, and the absence of a prime-degree factor of the
specified large size at every place. These are the four tests of
the finite type. We give the counting argument in the next
paragraph.\label{cor:paired-good-share-4}

For the two root tests, let \(g\) be the chosen level prime, let \(r\) be the rank, and put \(N=r\). For a degree-\(j\) label, the order of its multiplicative group and its \(\ell\)-primary part are
\[
  A_j=g^j-1,\qquad A_{j,\ell}=\ell^{v_\ell(g^j-1)}.
\]
 The
squarefree mass of labels that contain specified distinct
irreducibles of degrees $j_1,\ldots,j_u$ is at most
$\prod_i A_{j_i}^{-1}$.

Strip off the specified factors and sum
the remaining weights, which have total at most one in the
remaining degree. This last assertion also follows without the
character formula: if $I_j$ is the number of irreducibles of
degree $j$ with nonzero constant term, then $jI_j\le A_j$ by
counting their distinct roots in $\mathbb F_{g^j}^\times$.
Thus the mass of polynomials with factor counts $m_j$ is at
most $\prod_j(j^{m_j}m_j!)^{-1}$, the probability of that cycle
count under a uniform permutation. This proves domination for
any event specified by the factor degrees.

\begin{lemma}[Mass of the four excluded finite types]
\label{lem:finite-type-event-bounds}
Use the squarefree subprobability measure and the hypotheses of
\eqref{eq:row-bad-types}. Write \(m_{\rm pair},m_{\rm scal},m_{\rm div},m_{\rm cyc}\)
for the probabilities of a repeated prime-to-\(\ell\) root part, a match
with a prescribed scalar or its inverse, a common proper divisor degree,
and absence of the prescribed large prime degree, respectively. Then
\[
 \begin{aligned}
 m_{\rm pair}&\le 23f\log(g)/\sqrt g,&
 m_{\rm scal}&\le2fm\left(\frac1{4\min(q,d')}+\frac{1+\log r}{g-1}\right),\\
 m_{\rm div}&\le\tfrac r2(7/8)^f,&
 m_{\rm cyc}&\le\exp\!\left(-\frac f{5\log r}\right).
 \end{aligned}
\]
The pair test ranges over the nonexcluded \(\ell\), and the last two
bounds use the product over the \(f\) residue factors. The separate
summand \(3f/g\) in \eqref{eq:row-bad-types} allows for the missing
squarefree mass.
\end{lemma}

\begin{proof}
\proofstep{Repeated roots in the squarefree mass.}
When two different roots of a single irreducible of degree $j$
have the same prime-to-$\ell$ part, they are $z,z^{g^a}$ with
$1\le a<j$. Their equality requires $z^{g^a-1}$ to lie in the
$\ell$-primary group. In the cyclic group of order $A_j$, this
allows at most $(g^{\gcd(a,j)}-1)A_{j,\ell}$ choices of $z$.
For a nonexcluded prime, $A_{j,\ell}\le g^{j/4}$ by its
exclusion test.

Since \(\gcd(a,j)\le j/2\), the root count is bounded by
\[
  (g^{\gcd(a,j)}-1)A_{j,\ell}\le g^{j/2}g^{j/4}=g^{3j/4}.
\]
 An equality between distinct roots can occur only
for $\ell\mid A_j$. There are at most $j\log_2 g$ such primes.
Every represented irreducible has at least $j$ winning roots,
because the test is invariant under Frobenius. Dividing the root
count by $j$ and then by $A_j$ bounds the mass for one factor by
\[
 \frac{4096}{4095}\log_2 g
       \sum_{j\ge2}(j-1)u^j,
 \qquad u=g^{-1/4}\le1/8.
\]

For a pair of distinct irreducibles of degrees $i,j$ use the
cyclic groups on the two roots. The common prime-to-$\ell$
part lies in the group of order $g^{\gcd(i,j)}-1$, and the two
primary parts give at most $A_{i,\ell}A_{j,\ell}$ possibilities.
Thus the number of ordered root pairs is at most $g^{3(i+j)/4}$.
The union over the prime factors of $A_i A_j$, followed by
stripping the two factors from the label, bounds the mass by
\[
 \left(\frac{4096}{4095}\right)^2\log_2 g
                \sum_{i,j\ge1}(i+j)u^{i+j}.
\]
The sums are $u^2/(1-u)^2$ and $2u^2/(1-u)^3$. Together they
fit $23\log g/\sqrt g$, with a constant independent of $N$.
This also bounds the whole-place failure by the union over places.
\label{cor:paired-type-mass-1}

\proofstep{A fixed scalar and its root count.}
For a fixed scalar in $\mathbb F_g^\times$, the number of
roots in $\mathbb F_{g^j}^\times$ with matching prime-to-$\ell$
part is at most $A_{j,\ell}$. It suffices to take the union over
primes dividing $A_j$. If a match occurs at another prime the
root is the scalar itself, which is already counted. Divide by
$j$ for the number of roots of each irreducible, and by $A_j$
for the mass of labels containing it.
For fixed \(j,\ell\), these two divisions bound the mass from a fixed scalar by
\[
  \frac{1}{A_j}\,\frac{A_{j,\ell}}{j} .
\]
 At $j=1$, the sum of the
primary parts of $g-1$, divided by $g-1$, is at most
$3/(16\min(q,d'))$. Indeed the three parts for $2,q,d'$ have
size at least $16,q,d'$, and the product of any two is at least
$16\min(q,d')$.

For previous sum \(S\ge2\) and product \(P\), adjoining a primary part
\(c\ge2\) gives
\[
 \frac{S+c}{cP}\le\frac S P,
 \qquad cS-(S+c)=(c-1)(S-1)-1\ge0 .
\] At $j\ge2$, the product of the parts of
$A_j$ supported on $g-1$ is at most $j(g-1)$ by the lifting
formula for powers, also at two since $g\equiv1\pmod4$.
The product of the other parts is at most $A_j/(g-1)$.
In each group the sum of its primary parts is at most the
product. Thus the degree $j$ contribution is at most
$(g-1)/A_j+1/(j(g-1))$. Summing the resulting estimate uses both
\[
  \frac{g-1}{g^j-1}=\frac{1}{1+g+\cdots+g^{j-1}}\le g^{1-j},\qquad \sum_{j=2}^N\frac1j\le\log N.
\]
 This proves the scalar term
in \eqref{eq:row-bad-types} after summing over the $2fm$
scalars and inverse scalars.
\label{cor:paired-type-mass-2}

\proofstep{The divisor-degree and cycle tests.}
A divisor of degree $a$ for a squarefree label corresponds to
an invariant $a$-subset of a permutation in the degree-count
bound. For $0<a<N$, count stable $a$-subsets with weight eight,
except give weight four when the permutation fixes a point on
each side. For a fixed subset \(A\), let \(E_A\) be the event that
it is stable and \(F_A\) the event of a fixed point on
each side. Given \(E_A\), the two restrictions are
independent uniform permutations. The derangement formula
gives probability at least \(1/2\) of a fixed point on
each side, including a side of size one. Thus the mean
contribution from this subset is
\[
 \mathbb P(E_A)(8-4\mathbb P(F_A\mid E_A))\le7\mathbb P(E_A).
\]

When some subset is
stable the total weight is at least eight. In the weight-four
case, swap a fixed point on each side to obtain a different
stable subset of the same size. Summing the stabilization
probability over subsets gives one. Hence the chance of a
stable $a$-subset is at most $7/8$. Independence of the places
and a union over $1\le a\le N/2$ give the fourth term of
\eqref{eq:row-bad-types}.

A permutation can have at most one cycle of prime length in
$(N/2,N-3]$. For each such \(p\), the choice of its support, its cyclic ordering and the permutation on the complement gives
\[
  \Pr(\text{a cycle of length }p)=\binom Np\frac{(p-1)!(N-p)!}{N!}=\frac1p.
\]
 The sum of these reciprocals is at least
$1/(5\log N)$ for $N\ge100$. For $100\le N<16900$ use the sieve to count primes through
$16899$, and sum the lower bounds $10^{-12}\lfloor10^{12}/p\rfloor$.
Multiplying by $5\lfloor\log_2 N\rfloor\cdot0.6931471803$ gives
at least one for every such $N$, using only integer arithmetic.

For larger
$N$, write $n_1=\lfloor N/2\rfloor$, $s_1=\lfloor\sqrt N\rfloor$
and $t=\#\{p:N/2<p\le N-3\}$. The central binomial lower
bound $4^{n_1}/n_1$ and its prime factorization give
\[
 \frac{N-1}{6}\log 4 \le
                 \big(\pi(\sqrt{2n_1})+t+4\big)\log N.
\]
Indeed the primes at most $\sqrt{2n_1}$ cost at most $2n_1$
each, those between $\sqrt{2n_1}$ and $2n_1/3$ cost at most
the primorial up to $2n_1/3$, bounded by $4^{\lfloor2n_1/3\rfloor}$,
and those between $2n_1/3$ and $n_1$ do not divide the binomial.
At most three primes near the upper endpoint were omitted from
the list for $t$.

Now $\pi(\sqrt{2n_1})\le(s_1+9)/3$ by counting
classes coprime to $6$, and $\log N\le(s_1+1)/64+7.70407$ by
$\log u\le u-1$ applied to $u=(s_1+1)/128$. With $s_1\ge130$,
expanding the two sides gives
\[
 (s_1+21)((s_1+1)/64+7.70407)+0.693148\le0.093147s_1^2.
\]
Since $\log 2>0.69314718$, substitution gives $5t\log N\ge N$.
Each reciprocal in the sum is at least $1/N$. The cycle-count domination gives
$1-1/(5\log N)$ as an upper bound for the single-place mass
without one of the prescribed factors. Raising to the $f$th
power bounds the final event in \eqref{eq:row-bad-types}.
\label{cor:paired-type-mass-3}
\end{proof}


\begin{proposition}[A uniform bound on the removed rank]
\label{prop:paired-good-rank}
Use either run of Proposition~\ref{prop:linked-comparison-rows}
and its retained projector in
Proposition~\ref{prop:equivariant-selected-row}. Assume the
finite labels of Lemma~\ref{lem:finite-type-event-bounds} and
size bounds \eqref{eq:paired-bad-type-sizes}. Writing \(d_{\rm rem}\) for the removed rank at the lift, one has
\[
 \delta=\exp(-x^{e_z/7}),\qquad d_{\rm rem}\le\delta d_0,\qquad (1-\delta)d_0\le D_0.
\]
 The finite type bounds are applied
with the same parameters as the central projector trace.
\end{proposition}

\begin{proof}
\proofstep{Summing the finite tests over the selected labels.}

The size bounds for the type mass are
\begin{equation}\label{eq:paired-bad-type-sizes}
 f\ge x^2,\qquad f,m,r\le x^5,\qquad
 g\ge \exp(8x^\epsilon),\qquad
 \min(q,d')\ge\exp(x^{e_z/6}),\qquad
 0<e_z\le\epsilon<1 .
\end{equation}
The exponents \(e_z,\epsilon\) are fixed for the family. On these bounds every term in \eqref{eq:row-bad-types},
including those with $\log g$, is $o(\exp(-x^{e_z/7}))$, uniformly
in the place, lift and selected point. For the logarithm term use
$\log g/\sqrt g\le O(g^{-1/4})$. For the last two terms use
$\log r\le5\log x$ and $f\ge x^2$. At each level prime the resulting family bound is
\[
  4\beta_k\le\delta,\qquad \delta=\exp(-x^{e_z/7})\le\tfrac14.
\]
This threshold is common to the two runs and all the readings.
By \eqref{eq:row-bad-share} the removed rank is at most
$\delta d_0$. Using the tied rank equality from
Section~\ref{subsec:selected-cusp-reading} at each lift gives
$(1-\delta)d_0\le D_0$ on the retained row. Its numerator
comparisons use the corresponding eigenspaces in the same row.
\label{cor:paired-good-share-5}
\end{proof}
